\section{Experiments}
\label{sec:experiments}

\subsection{Protocols and Quantities}

Two sources of decision states are used. The controlled environment generates
episodes with a specified fraction of harmful candidates, so the harmful
fraction can be raised until routing is hazardous. The deployment states are
the saved decision traces of a response-aware router on six standard traffic
benchmarks: for each dataset, 32 fixed target sensors, four sequential
decision steps per episode, a 16-candidate pool, a context budget of four, and
a shortlist of four candidates. Each state carries the candidate scores, the
compact response summaries, and the true marginal utility of every
shortlisted candidate, which makes the harmful event observable for
evaluation while remaining invisible to the router.

Calibration and evaluation are separated by target sensors: the threshold is
calibrated on half of the targets and evaluated on the other half. We report
the harmful-selection mass (expected harmful selections per decision), the
coverage (fraction of decisions on which the router acts), the harmful share
among accepted decisions, and the mean true utility of accepted selections.
Intervals are target-level bootstrap intervals with 1000 replicates.

\subsection{Controlled harm ladder}

The controlled environment raises the fraction of harmful candidates while
keeping everything else fixed, so the effect of hazard on both gates is
visible without any change of protocol. Table~\ref{tab:harm-ladder} reports the
result over five runs per level. Two effects stand out. First, the
uncalibrated gate becomes steadily more damaging: its negative-transfer rate
grows from .032 at a harmful fraction of .25 to .227 at .75, while the
calibrated gate stays between .0001 and .0012. Second, the price of that safety
is coverage: the calibrated gate accepts between .14 and .31 contexts per
episode, against roughly four for the uncalibrated router, so the budget is
paid in abstentions rather than in harm.

The ladder also shows where structural headroom comes from. At a harmful
fraction of .25 the sequential oracle improves on the standalone oracle by
$H_{\mathrm{state}}=.151$; at .50 the advantage halves to .075, and at .75 it
disappears entirely. When nearly every candidate is harmful, no ordering of
the candidates helps, and the value of state conditioning vanishes with it.
Structural headroom is therefore a property of a task that still contains
useful contexts, not of state conditioning as such.

% Generated by scripts/build_risk_harm_ladder.py -- do not edit by hand.
% source: results/raw/R200_harm_ladder_*/h*/result.json
\begin{table}[t]
\centering
\small
\caption{Controlled harm ladder over five runs per harmful fraction. The
calibrated gate uses the split-conformal lower bound at level $.05$; the
uncalibrated router acts whenever the estimated utility is positive. Gains are
loss reductions relative to the anchor-only predictor.}
\label{tab:harm-ladder}
\resizebox{\linewidth}{!}{%
\begin{tabular}{lrrrrrrr}
\toprule
Harmful fraction & Cal. gain & Cal. harm & Cal. contexts & Uncal. gain & Uncal. harm & Static gain & Oracle gain\\
\midrule
0.25 & .0722 & .0001 & 0.3 & .2734 & .0317 & .2752 & .4064 \\
0.5 & .0459 & .0006 & 0.2 & .2219 & .0854 & .2279 & .3379 \\
0.75 & .0365 & .0012 & 0.1 & .1179 & .2267 & .1251 & .2130 \\
\bottomrule
\end{tabular}}
\end{table}


\subsection{A budget on harmful selections}

Table~\ref{tab:budget} reports the risk-controlled gate on deployment states.
At a budget of two harmful selections per hundred decisions, the gate retains
5.3\% to 10.7\% of decisions, and the realised mass tracks the nominal budget
within the bootstrap intervals on every dataset. Raising the budget to five
per hundred roughly doubles the retained coverage. The gate therefore converts
a safety requirement into a coverage number, which is the quantity an operator
can plan with.

% Generated by scripts/build_risk_tables.py -- do not edit by hand.
% source: results/derived/R200_risk_control/summary.json (P5)
\begin{table}[t]
\centering
\small
\caption{Risk-controlled gate on deployment states. Each cell reports the
realised harmful-selection mass and the retained fraction of decisions
(coverage) at the nominal budget $\alpha$. The threshold is calibrated on
half of the target sensors and evaluated on the other half.}
\label{tab:budget}
\resizebox{\linewidth}{!}{%
\begin{tabular}{lrrrr}
\toprule
Dataset & $\alpha=0.005$ & $\alpha=0.01$ & $\alpha=0.02$ & $\alpha=0.05$ \\
\midrule
METR-LA & .0067 / .022 & .0145 / .042 & .0252 / .065 & .0565 / .130 \\
PEMS-BAY & .0050 / .018 & .0097 / .031 & .0183 / .053 & .0462 / .113 \\
PEMS03 & .0064 / .029 & .0111 / .050 & .0205 / .089 & .0533 / .183 \\
PEMS04 & .0035 / .045 & .0085 / .067 & .0183 / .107 & .0467 / .200 \\
PEMS07 & .0070 / .034 & .0130 / .058 & .0260 / .099 & .0576 / .186 \\
PEMS08 & .0046 / .034 & .0076 / .052 & .0155 / .088 & .0506 / .195 \\
\bottomrule
\end{tabular}}
\end{table}


The uncalibrated gates, by contrast, accept far more and pay for it. Acting
whenever the score is positive accepts 45.0\% of decisions, of which 39.6\% are
harmful, for a harmful-selection mass of .179. Always selecting the top
candidate raises the mass to .495 and lowers the mean utility of accepted
selections to .0012, because most candidate states have a negative marginal.
The calibrated gate at $\alpha=.02$ keeps 86\% of the mean accepted utility of
the sign gate while removing 87\% of its harmful selections.

\subsection{Certification is not control}

Table~\ref{tab:certification} compares the two ways of using the same score.
The interval gate abstains unless the conformal lower bound of a candidate is
positive; the budget gate abstains unless the calibrated threshold is reached.
At the same nominal level, the interval gate accepts 1.0\% to 2.2\% of
decisions, while budget control accepts 5.3\% to 10.7\%, a factor of 4.3 to
9.1 more coverage, at a realised harmful mass at or below the budget. The
reason is the one identified in Section~\ref{sec:analysis}: the conformal
radius is a quantile of the estimation error, and in these states the error is
several times larger than the utility it estimates, so the lower bound is
negative almost everywhere.

% Generated by scripts/build_risk_tables.py -- do not edit by hand.
% source: results/derived/R200_risk_control/summary.json (P5, alpha = .02)
\begin{table}[t]
\centering
\small
\caption{Interval certification against budget control at the same nominal
level $\alpha=.02$. The interval gate accepts a decision only when the
conformal lower bound is positive; the budget gate accepts when the score
reaches the calibrated threshold. The last column is the coverage ratio of
budget control to interval certification.}
\label{tab:certification}
\resizebox{\linewidth}{!}{%
\begin{tabular}{lrrrrr}
\toprule
Dataset & Interval coverage & Budget coverage & Harm mass & Gain & Ratio\\
\midrule
METR-LA & .015 & .065 & .0252 & .0001 & 4.3$\times$ \\
PEMS-BAY & .012 & .053 & .0183 & .0068 & 4.3$\times$ \\
PEMS03 & .010 & .089 & .0205 & .0020 & 8.9$\times$ \\
PEMS04 & .022 & .107 & .0183 & .0061 & 4.9$\times$ \\
PEMS07 & .011 & .099 & .0260 & .0020 & 9.1$\times$ \\
PEMS08 & .014 & .088 & .0155 & .0028 & 6.5$\times$ \\
\bottomrule
\end{tabular}}
\end{table}


\subsection{Transfer of a calibrated threshold}

Table~\ref{tab:transfer} calibrates the threshold on five datasets and applies
it to the sixth. Realised harm stays within .0155--.0237 against a nominal
budget of .02, so the shift margin of Proposition~\ref{prop:shift} is small
across these benchmarks. The margin is not free: at a looser budget of .05 the
same protocol overshoots on one dataset (.070 against .050), which shows that
transfer quality degrades as the accepted set grows. Calibrating across
independent router training runs behaves the same way.

% Generated by scripts/build_risk_tables.py -- do not edit by hand.
% source: results/derived/R200_risk_control/summary.json (P6)
\begin{table}[t]
\centering
\small
\caption{Threshold transfer. The threshold is calibrated on every other
dataset at $\alpha=.02$ and evaluated on the held-out dataset. Realised harm
stays close to the nominal budget, and the last column reports the harmful
share among accepted decisions.}
\label{tab:transfer}
\resizebox{\linewidth}{!}{%
\begin{tabular}{lrrrr}
\toprule
Held-out dataset & Threshold & Harm mass & Coverage & Harmful share\\
\midrule
METR-LA & .360 & .0201 & .053 & .377 \\
PEMS-BAY & .355 & .0222 & .061 & .366 \\
PEMS03 & .353 & .0237 & .093 & .255 \\
PEMS04 & .370 & .0155 & .091 & .171 \\
PEMS07 & .363 & .0183 & .081 & .226 \\
PEMS08 & .359 & .0205 & .098 & .209 \\
\bottomrule
\end{tabular}}
\end{table}


\subsection{What the score can and cannot certify}

A mass budget bounds how many harmful selections the router makes. A share
budget, which bounds how often an accepted decision is harmful, is a different
requirement, and the same score does not support both. Table~\ref{tab:frontier}
reports the risk--coverage frontier: the smallest harmful share that any
threshold on the response-aware score can achieve at a given coverage. The
frontier is high and flat. At a coverage of 2\% the best achievable harmful
share is .202 on average and exceeds .24 on three datasets; at full coverage it
reaches .495. Consequently a share budget of 10\% is infeasible on five of the
six datasets at any coverage above 2\%, and a calibration procedure that
targets it returns no threshold at all.

% Generated by scripts/build_risk_tables.py -- do not edit by hand.
% source: results/derived/R200_risk_control/summary.json (P7)
\begin{table}[t]
\centering
\small
\caption{Risk--coverage frontier of the response-aware score. Each entry is the
smallest harmful share that any threshold on this score can achieve at the
stated coverage, computed on the evaluation targets. The frontier bounds every
threshold rule, including the calibrated gate.}
\label{tab:frontier}
\resizebox{\linewidth}{!}{%
\begin{tabular}{lrrrrrrr}
\toprule
Dataset & 0.02 & 0.05 & 0.10 & 0.20 & 0.30 & 0.50 & 1.00\\
\midrule
METR-LA & .308 & .363 & .418 & .456 & .469 & .485 & .514 \\
PEMS-BAY & .287 & .351 & .403 & .417 & .428 & .466 & .528 \\
PEMS03 & .241 & .221 & .235 & .302 & .341 & .388 & .497 \\
PEMS04 & .048 & .088 & .164 & .233 & .274 & .335 & .477 \\
PEMS07 & .205 & .221 & .262 & .319 & .353 & .400 & .494 \\
PEMS08 & .122 & .146 & .184 & .263 & .306 & .365 & .459 \\
\midrule
Mean & .202 & .232 & .278 & .332 & .362 & .406 & .495 \\
\bottomrule
\end{tabular}}
\end{table}


The frontier also explains why mass control works. Reducing the accepted set
removes harmful selections in proportion to their mass, but it does not
concentrate the remaining decisions on candidates that are reliably useful:
the score orders candidates well enough to lower the count of harmful
selections, not well enough to certify a low harmful rate among the decisions
it keeps. Improving the precision of the top of the ranking, rather than the
average ranking quality, is therefore the concrete target for a utility model
that is meant to support a share guarantee.

% Generated by scripts/build_risk_tables.py -- do not edit by hand.
% source: results/derived/R200_risk_control/summary.json (P5)
\begin{table}[t]
\centering
\small
\caption{Operating points. Each cell reports coverage and the harmful share
among accepted decisions. The budget gate is calibrated at $\alpha=.02$; the
final column reports its harmful-selection mass and the mean true utility of
its accepted selections.}
\label{tab:objectives}
\resizebox{\linewidth}{!}{%
\begin{tabular}{lrrrrr}
\toprule
Dataset & Always & Sign gate & Interval gate & Budget gate & Mass / gain\\
\midrule
METR-LA & 1.00 / .514 & .47 / .484 & .015 / .298 & .065 / .389 & .0252 / .0001 \\
PEMS-BAY & 1.00 / .528 & .47 / .460 & .012 / .250 & .053 / .349 & .0183 / .0068 \\
PEMS03 & 1.00 / .497 & .48 / .384 & .010 / .257 & .089 / .230 & .0205 / .0020 \\
PEMS04 & 1.00 / .477 & .43 / .315 & .022 / .055 & .107 / .171 & .0183 / .0061 \\
PEMS07 & 1.00 / .494 & .43 / .388 & .011 / .206 & .099 / .261 & .0260 / .0020 \\
PEMS08 & 1.00 / .459 & .42 / .347 & .014 / .108 & .088 / .176 & .0155 / .0028 \\
\bottomrule
\end{tabular}}
\end{table}


\subsection{End-to-end selective routing}

% TODO(R201): the end-to-end run applies the calibrated gate inside the
% sequential router and reports episode-level gain, negative-transfer rate and
% the number of contexts used at each budget.

\subsection{State-conditional thresholds}

Partitioning the states into score quintiles and calibrating a threshold per
bin reduces coverage at a fixed budget: at $\alpha=.02$ the conditional gate
retains 5.8\% of decisions against 8.4\% for the global threshold, at the same
realised mass. With four thousand calibration states per bin, the additional
variance of the local estimates outweighs their adaptivity. We therefore keep
the global threshold and report conditional calibration as an ablation rather
than as a component of the method.
